\documentclass{article}

\usepackage[submission]{../../../colm/2026/colm2026_conference}
\usepackage{booktabs}
\usepackage{hyperref}
\usepackage{url}

\title{Decode Microkernel Consolidation for Hybrid Reasoner Serving}
\author{Anonymous authors\\Paper under double-blind review}

\begin{document}
\maketitle

\begin{abstract}
This draft reports the current state of Decode Microkernel Consolidation
(DMC), a positive-method pivot produced after a preregistered HybridKernel
boundary-fusion gate failed. DMC targets repeated launch-heavy decode
micro-operations rather than attention/SSM boundaries. The current evidence is
a preregistered Phase 1 replay gate: a trace-derived packed Triton replay path
reduces GPU-side replay latency while preserving numerical equivalence on
three primary, three same-family, and three cross-family schedules. This is not
yet an end-to-end serving speedup claim. The next preregistered gate must show
that the same method improves real vLLM/serving decode latency on frozen
prompt slices.
\end{abstract}

\section{Introduction}

Hybrid and mixture-of-experts language models expose many small decode-time GPU
operations. A failed HybridKernel profiler gate found no separable
attention/SSM boundary-local overhead, but the same server-side Nsight
artifacts showed dense repeated micro-operation families during decode. DMC
asks whether those repeated operations can be consolidated into fewer GPU
launches without changing outputs.

The submission status is deliberately conservative. Phase 1 is a replay
implementation gate, not a serving result. The paper is not camera-ready until
Phase 2 shows an end-to-end serving gain with frozen prompts, model outputs,
launch audits, and paired uncertainty.

\section{Related Work}

This interim draft does not make external novelty claims beyond the repository
audit. The camera-ready candidate must include a hand-verified citation audit
covering serving launch overhead, compiler fusion, CUDA Graphs, Triton kernels,
and hybrid/MoE decode systems. Until that audit is complete, the only cited
evidence is internal artifact evidence listed in the reviewer pack.

\section{Method}

DMC builds a fixed replay schedule from committed server-side Nsight Systems
\texttt{cuda\_gpu\_kern\_sum} logs in the killed HybridKernel packet. For each
source row, the runner maps kernel-family weights into a packed
gating/routing/state-update replay schedule. The baseline executes repeated
PyTorch CUDA micro-operations; the consolidated path executes one Triton packed
gating replay kernel. The checker verifies CUDA-event timing samples, launch
audit artifacts from \texttt{torch.profiler}, saved output tensors, source
artifact hashes, fixed packet identity, and numerical equivalence.

The claim boundary is strict: DMC Phase 1 validates a replay mechanism. It does
not validate model quality, vLLM integration, attention/SSM boundary fusion, or
end-to-end serving acceleration.

\section{Experiments}
\label{sec:experiments}

Phase 1 used the fixed PASS packet
\url{experimental/decode_microkernel/phase1/results/dmc_phase1_20260508T000525Z}.
The gate was preregistered in
\url{experimental/decode_microkernel/phase1/preregister_dmc_phase1.md}. The
checker output is
\url{experimental/decode_microkernel/phase1/results/dmc_phase1_20260508T000525Z/checker_result.json}.

Each row used at least 100 warmup iterations and 1000 measured CUDA-event
iterations. All nine rows were checked against saved output tensors and
per-row launch-audit artifacts.

\section{Results}

\begin{table}[t]
\centering
\scriptsize
\begin{tabular}{lrrrr}
\toprule
Role & Rows & Median latency reduction & 95\% CI low & Max abs error \\
\midrule
Primary hybrid & 3 & 0.9650953811812909 & 0.9618806519742926 & 7.152557373046875e-07 \\
Same-family control & 3 & 0.9657258621103347 & 0.9656092336667471 & 5.960464477539062e-07 \\
Cross-family falsification & 3 & 0.9655504614632344 & 0.9654552838317912 & 5.960464477539062e-07 \\
\bottomrule
\end{tabular}
\caption{Phase 1 replay-gate metrics from the checker output cited in
Section~\ref{sec:experiments}.
The preregistered thresholds were median latency reduction at least 0.08 for
primary and same-family roles, at least 0.05 for cross-family roles, bootstrap
95\% CI lower bound greater than 0 for every role, launch reduction at least
0.25 in every row, max absolute error at most 1e-2, and max relative error at
most 1e-2.}
\end{table}

The checker returned
\texttt{PASS\_DMC\_PHASE1\_CONSOLIDATED\_REPLAY}. The minimum row-level launch
reduction was 0.9916666666666667. The maximum relative error was
8.915998250813573e-07.

\section{Discussion}

The Phase 1 result is a useful positive signal because it shows that a
trace-conditioned packed replay can consolidate launch-heavy decode
micro-operations with exact artifact checks. It is not sufficient for a COLM
method paper by itself. The replay baseline intentionally exposes launch
overhead; an optimized serving stack may already remove some of that overhead.
The next gate therefore measures real serving rows and requires generated-token
matches, end-to-end decode latency reduction, tokens/sec gain, and launch
audits.

\section{Limitations}

Phase 1 has only nine replay rows and three rows per role. Bootstrap intervals
over three rows should not be treated as paper-strength uncertainty. The
workload is trace-derived but synthetic; it does not exercise full model
generation, KV-cache behavior, batching, sampling, prefill/decode interaction,
or existing vLLM compiler and CUDA Graph optimizations. Same-family evidence
includes one row whose Phase 0 SQLite summary was empty, although Phase 1
derives that schedule from the committed server-side Nsight log and records
that provenance explicitly. Phase 2 must resolve these limitations before
camera-ready status.

\section{Reproducibility Statement}

The current Phase 1 packet was produced on commit
\texttt{86fd5b027ea6259cbba9be60a81427efa2d525f5} on an NVIDIA RTX PRO 6000
Blackwell Server Edition GPU. Runtime was less than 0.01 GPU-hours for the
canonical replay packet. Reproduction uses:

\begin{verbatim}
HF_HOME=/workspace/hf_cache TRITON_CACHE_DIR=/workspace/.triton_cache \
./.venv_gpu/bin/python experimental/decode_microkernel/phase1/run_dmc_phase1.py \
  --run-id dmc_phase1_20260508T000525Z

./.venv_gpu/bin/python experimental/decode_microkernel/phase1/check_dmc_phase1.py \
  experimental/decode_microkernel/phase1/results/dmc_phase1_20260508T000525Z
\end{verbatim}

All claims in this draft cite
\url{experimental/decode_microkernel/phase1/results/dmc_phase1_20260508T000525Z}.

\section{Ethics Statement}

This work is a systems optimization study. The current artifacts do not involve
human-subject data, user data, or deployment claims. The main ethical risk is
overclaiming a replay microbenchmark as a serving speedup; the paper explicitly
forbids that interpretation until Phase 2 passes.

\end{document}
